\documentclass[11pt]{article}

\usepackage[utf8]{inputenc}
\usepackage[T1]{fontenc}
\usepackage{textcomp}
\usepackage[a4paper,margin=1in]{geometry}
\usepackage{parskip}
\usepackage{hyperref}

% Reviewer/editor comments are set in italics; responses follow in normal type.
\newenvironment{italic}{\itshape}{\par}

\title{Response to the Editor and Reviewers\\[0.5em]
\large ITMO Regulations to Ratchet up Ambition}
\author{}
\date{}

\begin{document}

\maketitle

\noindent We thank the Editor and both Reviewers for their careful and constructive assessment of the manuscript. % Below we reproduce each comment in italics and provide our response underneath. Changes are highlighted in the tracked-changes version of the manuscript.

\bigskip

\section*{Editor}

\begin{italic}
Dear Dr.\ Fabre,

The editors of \emph{Global Environmental Change} have secured two reviews of your manuscript. Both reviewers are positive about the paper and signal its strengths and resonance. However, both also highlight a series of minor revisions to be made. Both reviewers have provided a thorough, constructive set of comments that are self-explanatory. Please also pay particular attention to Reviewer~1's request for changes to be made to the manuscript's structure. Please find their reviews below.

We are happy to invite you to do further work on your manuscript, attending carefully to each of the concerns raised, before re-submitting it for our further consideration. In this case, we will likely approach one or both reviewers for their further opinion.

In your re-submitted version, please provide an accompanying letter explaining in detail the responses made to each point by the reviewers and a version of your manuscript with tracked changes included in addition to a clean version. Please also consult the guide for authors carefully to ensure your manuscript is formatted correctly (e.g.\ Harvard referencing, avoiding or minimising the use of footnotes, etc.) to avoid additional work at the resubmission stage.
\end{italic}



We thank the Editor.



\section*{Reviewer 1}

\begin{italic}
Dear Author, thank you for the opportunity to review the Perspective article entitled ``ITMO Regulations to Ratchet up Ambition''. The topic is quite relevant and with important timing, as some countries are still submitting/revising their NDCs, and this could improve next NDC submission cycles. I have, however, identified some minor aspects that could be improved before the manuscript can be considered for publication. Below, I provide my detailed comments, divided into major and minor points.
\end{italic}



We thank the Reviewer for the positive assessment.



\subsection*{Major comments}

\begin{italic}
As a major comment, but a fairly easy one, I would suggest that the author make some changes to the overall structure of the Perspective, merging sections or changing some sections to subsections (e.g.: sections 4, 5 and 6 being merged as one section with subsections).
\end{italic}






\subsection*{Minor comments}

\begin{italic}
Page 4, line 96, footnote 3: The author uses Angola as an example of a country that has not ratcheted up its ambition. However, one could also argue that the Angolan government is trying to pledge what it considers feasible after reassessing the extent to which the country has been able to reduce emissions and approach its 2025 target. I would suggest keeping the wording as country-neutral as possible, in order to strengthen support for the author's proposal rather than emphasising perceived shortcomings in specific national contexts.
\end{italic}






\begin{italic}
Page 4, line 125, ``Until now, I conflated conditional and unconditional NDC targets.'': I dare say that the author would make a stronger argument if this part, and the sequential paragraph, appeared before, as this can be a decisive argument to Global South countries.
\end{italic}






\begin{italic}
Page 6, line 198, ``usewould'', typo.
\end{italic}






\begin{italic}
Page 7, line 208, ``2\,\textdegree C with a 75\% probability'' \& footnote 6, could use a citation.
\end{italic}






\begin{italic}
In all subsection titles, I suggest that the author take out the dot at the end.
\end{italic}






\begin{italic}
At this stage, I recommend minor revisions for your Perspective paper. I hope the comments above will help strengthen the manuscript.
\end{italic}



We thank the Reviewer.



\section*{Reviewer 2}

\begin{italic}
The study and its methodology are transparent, well-articulated, and present a commendable proposal for a way forward. I found the study to be highly interesting and informative. However, I would appreciate it if you could clarify the assumptions underlying your proposal regarding the method of defining the ambition of the Nationally Determined Contributions (NDCs) and the grouping of the coalition.
\end{italic}



We thank the Reviewer for the positive assessment.



\begin{italic}
While the general approach is insightful, I think certain aspects of the methodology could be more clearly defined, particularly regarding how the Business-As-Usual (BAU) scenario is delineated. A clearer definition of BAU will enhance the understanding of the proposed methodology and the implications for future analyses.
\end{italic}






\subsection*{Main comments}

\begin{italic}
\textbf{1. Assumptions on BAU:} It would be beneficial to outline the key assumptions related to the Business-As-Usual (BAU) scenario in greater detail. This clarification will help readers critically evaluate the validity of your approach. I propose utilizing a current policies projection instead of the BAU scenario, as BAU is often considered a no-policy baseline, which may appear unambitious (e.g., SSP 8.5) and typically does not account for existing policies. For guidelines on defining a current policies scenario, please refer to the UNEP Emissions Gap Reports (Chapter~3), as well as studies such as den Elzen et al.\ (2019) and the latest UNEP Emissions Gap Report 2025 (Olhof et al.\ 2025). These sources present methodologies for assessing current policy scenarios, which include national independent studies, global model studies, and governmental analyses. It states: ``This scenario projects global GHG emissions assuming all currently adopted and implemented policies as of a certain specific date are realized and that no additional measures are undertaken.'' I recommend including a current policies scenario or at least referencing it in your discussion. Additionally, please elaborate on the BAU scenario and specify the sources of information for current policies.
\end{italic}






\begin{italic}
\textbf{2. Literature Sources on Evaluating NDC Ambition and BAU (Section 2.3):} A brief description of the reference sources utilized in the study to analyze the ambition of NDCs would be beneficial. There are various methods to assess the ambition of NDCs, as discussed in the literature, such as comparing with least-cost scenarios or equity outcomes. More references that analyze this ambition would enhance your work.
\end{italic}






\begin{italic}
\textbf{3. Current Policies and Ambition:} It is crucial that a country adopts policies that match its ambition, as clearly explained in Nascimento et al.\ (2024). Both Nascimento et al.\ and the UNEP Emissions Gap Reports (Chapter~3) show that certain countries already have NDCs that lead to surpluses from the date of submission, such as India, Russia, and Turkey. The issue of surpluses is significant, as some countries have relatively unambitious NDCs.
\end{italic}


Explain here that as long as the sum of NDCs is below the Paris target, a country can have a surplus and still be considered unambitious. Carbon trading should equalize abatement costs across countries and NDCs should only reflect burden-sharing, not abatement activity.



\begin{italic}
\textbf{4. Section 2.3 -- Line 103:} While I appreciate your proposal, it's important to note that the concept of ``equal share'' is debated in climate policy discussions. I encourage you to elaborate further on what constitutes an ``equal share'' in this context. It would be valuable to reference existing literature that examines different interpretations and applications of this concept, as a more comprehensive understanding will strengthen your argument.
\end{italic}

after "under that 2°C scenario", add "(thereby spanning a wide range of candidate burden-sharing rules, \citealp{van_den_berg_implications_2020})"




\begin{italic}
\textbf{5. Line 149:} I respectfully disagree with the assertion made here. I believe that nearly all countries, including China, must be included in the scope of discussion if we are to meet the Paris Climate Agreement goals. It's essential to consider that China's per capita emissions are significantly higher than those of the European Union and the world average. The effectiveness of our global climate strategy hinges upon the collective ambition of these major emitters. Thus, I think the assessment of NDC ambition should reflect the realities of emissions across different countries to capture a more accurate picture of progress toward climate targets.
\end{italic}


added ``(including China), which should decarbonize faster and provide climate finance abroad'' after ``the ratcheting up of ambition should arguably be borne by industrialized countries''



\subsection*{Specific comments}

\begin{italic}
\textbf{P2, line 59:} This section is somewhat undefined. The Climate Action Tracker presents current policies, which I believe should be utilized here. An alternative would involve emissions projections based on equity considerations, but this approach is not suitable for this context.
\end{italic}

Add explicit reference to current policies




\begin{italic}
\textbf{P3, line 66:} For instance, current policies projections can be derived from a literature assessment, as demonstrated in Chapter~3 of the UNEP Emissions Gap Report (Olhoff et al., 2025) and in publications like den Elzen et al.\ (2019). However, the data is often limited at the country level for non-G20 economies, with the exception of downscaled results, such as those from the NGFS, which may not be completely reliable.
\end{italic}


add suggested citations in footnote



\begin{italic}
\textbf{Page 3, line 78:} Yes, the footnote references multiple studies. You could also mention Chapter~3 of the UNEP Emissions Gap Report, which provides an assessment of studies for UNEP, using harmonization for all studies in terms of Land Use, Land-Use Change, and Forestry (LULUCF). References to the UNEP Emissions Gap Report or den Elzen et al.\ (2019) would be appropriate. Refer to: Climate Analytics, 2026. \emph{1.5\,\textdegree C National Pathway Explorer}. Available at: \url{http://1p5ndc-pathways.climateanalytics.org}.
\end{italic}


do it



\begin{italic}
\textbf{Page 4, line 114, ambition gap:} This is not, for all readers, well defined. Please refer to the references, such as the UNEP emissions gap report, but also in literature it has been defined.
\end{italic}

add ``defined as the difference between the Paris target and the sum of NDCs'' and cite UNEP Emissions Gap Report (2025)




\begin{italic}
\textbf{Footnote 4:} This is not the best reference. I would focus on the UNEP Emissions Gap Report (Olhoff et al.\ 2025) or Nascimento et al.\ (2024), who reached similar conclusions regarding the current policies scenario. These are original works, whereas Van den Berg et al.\ primarily presented equity outcomes and an outdated BAU projection.
\end{italic}


add unep_emissions_2025 to van den berg citation 



\begin{italic}
\textbf{Line 159:} Reference multiple sources, including UNEP Emissions Gap Report (2025) and literature like Rogelj et al.\ (2023).
\end{italic}


do it



\end{document}
